
\chapter{Introdução}
\label{ch:intro}
O presente relatório foi desenvolvido no âmbito da unidade curricular de Sistemas Autónomos Inteligentes, do Mestrado em Engenharia Eletrónica Industrial e de Computadores da Universidade do Minho.

O trabalho centra-se no estudo e na implementação da cinemática do manipulador robótico myCobot~280, constituído por seis articulações rotativas. O objetivo consiste em determinar a posição do centro de preensão e a orientação da garra a partir das variáveis articulares, bem como calcular as configurações necessárias para alcançar diferentes posições no espaço de trabalho.

Para o desenvolvimento do projeto, foi utilizado o MATLAB, responsável pela implementação dos modelos cinemáticos e pelo controlo dos movimentos, e o CoppeliaSim, utilizado como ambiente de simulação para visualizar, testar e validar o comportamento do manipulador.

Ao longo do relatório são apresentados os fundamentos da cinemática direta e inversa, as respetivas formulações matemáticas, a implementação dos algoritmos em MATLAB e os resultados obtidos em ambiente de simulação.


\section{Fundamentos da cinemática}
\label{sec:fundamentos}

A cinemática de um manipulador robótico permite estabelecer a relação entre as variáveis articulares e a posição e orientação do \textit{end effector}. No caso do myCobot~280, constituído por seis articulações rotativas, esta análise divide-se em dois problemas fundamentais: a cinemática direta e a cinemática inversa.

A cinemática direta permite determinar a posição e orientação de um manipulador a partir dos ângulos das suas articulações. Por outro lado, a cinemática inversa permite determinar os ângulos articulares necessários para alcançar uma posição pretendida, considerando determinadas restrições de orientação.



\section{Enquadramento e objetivos}
\label{sec:enq-e-obj}

Os manipuladores robóticos desempenham um papel importante na automação industrial, sendo utilizados em tarefas que exigem precisão, repetibilidade e flexibilidade de movimentos. O seu funcionamento depende do conhecimento da geometria do manipulador e da relação entre as variáveis articulares e o posicionamento do \textit{end effector} no espaço de trabalho.

Neste contexto, o presente projeto centra-se no estudo do manipulador myCobot~280, recorrendo ao MATLAB para implementar os algoritmos cinemáticos e ao CoppeliaSim para simular e avaliar o comportamento do robô.

Os principais objetivos do trabalho são:

\begin{itemize}
    \item Desenvolver a formulação matemática da cinemática direta e inversa do manipulador;

    \item Implementar os modelos cinemáticos em MATLAB;

    \item Determinar a posição do centro de preensão e a direção da garra a partir dos ângulos articulares;

    \item Calcular as configurações articulares necessárias para alcançar posições pretendidas, considerando a inclinação da garra e os limites das juntas;

    \item Implementar movimentos no espaço das juntas e trajetórias no espaço cartesiano;

    \item Integrar o MATLAB com o CoppeliaSim para executar e visualizar os movimentos do manipulador;

    \item Explorar a utilização da visão artificial na identificação e localização de objetos;

    \item Validar os modelos implementados através da comparação entre os resultados calculados e os obtidos em simulação.
\end{itemize}

